\documentclass[10pt,a4paper]{amsart}
\usepackage[margin=19mm]{geometry}
\usepackage{amsmath,amssymb,mathtools}
\usepackage[scr=rsfs,cal=euler]{mathalfa}
\usepackage{xfrac}
\usepackage[unicode,hidelinks,bookmarks=false]{hyperref}
\newcommand{\ie}{i.e.\ }
\newcommand{\cf}{cf.\ }
\newcommand{\resp}{respectively\ }
\newcommand{\Subsection}{\subsection}
\providecommand{\nobreakdash}{\nobreak\hbox{-}}
\setlength{\emergencystretch}{2em}
\multlinegap0pt
\usepackage{amssymb}
\usepackage{xcolor}
\usepackage{algorithm}\usepackage{algpseudocode}
\usepackage{float}
\newtheorem{lemma}{Lemma}
\title{Fast and Stable Gradient Approximation for Bilinear Forms of Hermitian Matrix Functions}
\author{Navjot Singh, Kipton Barros, Xiaoye Sherry Li}
\date{Source: arXiv:2605.12801v1 (CC BY 4.0)}
\begin{document}
\maketitle

\noindent\emph{Source and licence.} Fast and Stable Gradient Approximation for Bilinear Forms of Hermitian Matrix Functions. Version arXiv:2605.12801v1 (CC BY 4.0), distributed by arXiv under the Creative Commons Attribution 4.0 International licence, http://creativecommons.org/licenses/by/4.0/, as recorded in the arXiv licence metadata for this exact version and shown on the abstract page https://arxiv.org/abs/2605.12801v1. Full licence notice: \texttt{licenses/arxiv-2605.12801-CC-BY-4.0.txt}.\par\smallskip

\noindent\textit{Editorial scope.} Counted unit: the article's lemma lem:app\_commutator (source0.tex lines 1172--1178). The article's complete original proof of it is delivered with it (source0.tex lines 1180--1278, one \texttt{proof} environment of 99 lines). No mathematics is written, repaired or re-derived: the statement and the proof are the article's own lines, in the article's own notation, and no retained line is substituted or re-flowed.  Retained uncounted as the source-local context the proof needs: lines 1018--1023.  Nothing was omitted from any retained span, so the package carries no omission record.  No retained line contains a citation, so the wrapper carries no bibliography and no bibliography entry.  No retained line contains a figure environment, an image-inclusion command or drawing code, so no image is delivered and the package declares no asset dependency.  Editorial formatting only, in addition: an AMS \texttt{amsart} preamble that restates the article's own declarations and macros used by the retained lines, namely the environments \texttt{lemma} on separate counters, together with the standard packages those lines need.  The retained environments print in this document as Lemma 1 in order of appearance, whereas the article prints the same items with the numbering of its own full document.  The complete native source of the article as distributed in the arXiv e-print for this exact version is bundled unchanged as \texttt{sources/200453/source0.tex}.
\begin{equation}
    G
    =
    \|u\|^2 Q\bigl((cc^\top)\circ F\bigr)Q^\top .
    \label{eq:app_G}
\end{equation}

\begin{lemma}[Vanishing commutator]
\label{lem:app_commutator}
If \(S^\top=-S\) and \(Se_1=0\), then
\[
    \operatorname{tr}(G^\top[T,S])=0 .
\]
\end{lemma}

\begin{proof}
Let
\[
    \widetilde S = Q^\top S Q .
\]
Since \(S^\top=-S\), we also have \(\widetilde S^\top=-\widetilde S\). Moreover,
using \(c=Q^\top e_1\) and \(Se_1=0\),
\[
    \widetilde S c
    =
    Q^\top S Q Q^\top e_1
    =
    Q^\top S e_1
    =
    0 .
\]
Because \(T=Q\Lambda Q^\top\),
\[
    Q^\top[T,S]Q
    =
    Q^\top(TS-ST)Q
    =
    \Lambda \widetilde S-\widetilde S\Lambda .
\]
Using the definition of \(G\) from \eqref{eq:app_G} and cyclic invariance of the trace,
\begin{align}
    \operatorname{tr}(G^\top[T,S])
    &=
    \|u\|^2
    \operatorname{tr}\!\left(
        \bigl((cc^\top)\circ F(\Lambda)\bigr)
        (\Lambda\widetilde S-\widetilde S\Lambda)
    \right).
\end{align}
It remains to show that the trace on the right-hand side is zero. Define
\[
    \mathcal{C}
    =
    \operatorname{tr}\!\left(
        \bigl((cc^\top)\circ F(\Lambda)\bigr)
        (\Lambda\widetilde S-\widetilde S\Lambda)
    \right).
\]
Expanding entrywise,
\begin{align}
    \mathcal{C}
    &=
    \sum_{i,j}
    c_i c_j F(\lambda_i,\lambda_j)
    (\lambda_j-\lambda_i)\widetilde S_{ji}.
    \label{eq:app_commutator_expansion}
\end{align}
The diagonal terms vanish because \(\lambda_j-\lambda_i=0\) when \(i=j\). For
\(i\neq j\), the divided-difference identity gives
\[
    F(\lambda_i,\lambda_j)(\lambda_j-\lambda_i)
    =
    f(\lambda_j)-f(\lambda_i).
\]
Hence
\begin{align}
    \mathcal{C}
    &=
    \sum_{i,j}
    c_i c_j
    \bigl(f(\lambda_j)-f(\lambda_i)\bigr)
    \widetilde S_{ji} \\
    &=
    \sum_{i,j}
    c_i c_j f(\lambda_j)\widetilde S_{ji}
    -
    \sum_{i,j}
    c_i c_j f(\lambda_i)\widetilde S_{ji}.
\end{align}
The two sums can be written in matrix form as
\[
    \mathcal{C}
    =
    (f(\Lambda)c)^\top \widetilde S c
    -
    c^\top \widetilde S f(\Lambda)c .
\]
Since \(\widetilde S\) is skew-symmetric,
\[
    c^\top \widetilde S f(\Lambda)c
    =
    - (f(\Lambda)c)^\top \widetilde S c .
\]
Therefore
\[
    \mathcal{C}
    =
    2(f(\Lambda)c)^\top \widetilde S c .
\]
Finally, \(\widetilde S c=0\), so \(\mathcal{C}=0\). Thus
\[
    \operatorname{tr}(G^\top[T,S])=0 .
\]
\end{proof}

\end{document}
